\scanpage{80}
sont pas nuls en général $=$ pour $\nu=1$. En effet,
\inserted{\unsure{prenons|supposons} pour $A'$ la duale de $A$.}
si \deleted{\unsure{$x$}} $x\in{}_{\ell^\nu}T(k)$\inserted{$\subset{}_{\ell^\nu}A(K)\subset{}_{\ell^\nu}A(\overline K)$} est non nul, il existe un
$x'\in{}_{\ell^\nu}A'(\overline K)$ tel que \deleted{\unsure{$\varphi$}} $\varphi(x,x')\ne0$, et \deleted{que} quitte
à agrandir $K$ en $K'$, on aura donc un
accouplement ${}_{\ell^\nu}A'(K)\times{}_{\ell^\nu}T(k)\longrightarrow{}_{\ell^\nu}k^*$ non nul $\ldots$.
\margin{79}
